\documentclass[11pt,a4paper]{article}
\usepackage[margin=25mm]{geometry}
\usepackage{amsmath,amssymb}
\usepackage[T1]{fontenc}
\usepackage[utf8]{inputenc}
\usepackage{lmodern}
\usepackage{enumitem}
\usepackage{needspace}
\usepackage{tikz}
\usepackage{xurl}
\usepackage[colorlinks=true,linkcolor=black,urlcolor=blue!55!black]{hyperref}
\hypersetup{pdfauthor={Luca Scharr},pdftitle={Tutorial 09 - Gaussian Splatting and NeRF}}
\setlength{\parindent}{0pt}
\setlength{\parskip}{0.45em}
\setlist{leftmargin=*,itemsep=0.15em,topsep=0.15em}
\emergencystretch=2em
\begin{document}
\begin{center}
{\Large\bfseries Tutorial 9. Gaussian Splatting and NeRF}\\[2mm]
Luca Scharr\\
Computer Vision 2026\\
University of Bern\\[1mm]
Interactive lab: \href{https://sosekie.github.io/Unibe-CVG-CV-Interactive-Labs/tutorial-09/}{\nolinkurl{sosekie.github.io/Unibe-CVG-CV-Interactive-Labs/tutorial-09/}}\\
Source code: \href{https://github.com/Sosekie/Unibe-CVG-CV-Interactive-Labs}{\nolinkurl{github.com/Sosekie/Unibe-CVG-CV-Interactive-Labs}}
\end{center}

\section*{Learning objectives}
\begin{itemize}[nosep]
\item Relate Gaussian shape to covariance, PCA, and function approximation.
\item Derive projection, compositing, and fitting behaviour from the underlying models.
\item Compare Gaussian Splatting and NeRF and evaluate reconstruction evidence.
\end{itemize}
\noindent\textbf{Key concepts.} Least squares, covariance, gradient descent, alpha compositing, multi-view geometry, volume rendering, and neural radiance fields.

Explain your reasoning and show relevant derivations.
The interactive lab provides experiments for each exercise.
\Needspace{12\baselineskip}\section*{1. Approximate a function in 1D}
\noindent\textbf{Objective.} Distinguish amplitude fitting, basis geometry, and identifiability.
\Needspace{9\baselineskip}\subsection*{1.1 Identifiability - Linear Fitting}
You fit two 1D Gaussian basis functions with exactly the same mean and width.
Only their amplitudes are free.
The target samples are fixed.
\begin{enumerate}[label=(\alph*),nosep]
\item What can least squares determine about the two amplitudes?
Justify your answer.
\item Would adding a penalty $\lambda\|a\|^2$ make the amplitudes identifiable from the observations?
\end{enumerate}
\Needspace{9\baselineskip}\subsection*{1.2 Model Assumptions - Signed Weights}
After an unconstrained amplitude fit, a nonnegative target is approximated using both positive and negative Gaussian coefficients.
Small regions of the approximation fall below zero.
\begin{enumerate}[label=(\alph*),nosep]
\item Which interpretation is consistent with the model used in this lab?
Justify your answer.
\item Why can a low sample error coexist with large, unstable positive and negative coefficients?
\end{enumerate}
\Needspace{7\baselineskip}\noindent\textbf{Interactive experiment 1.}\begin{enumerate}[nosep]
\item Fit the smooth-bump target with N = 2 and N = 8.
Record sample RMSE and basis rank.
\item Make two kernels identical in mean and width.
Predict the rank and possible amplitude solutions before fitting.
\item Try the signed-oscillation target.
Inspect the individual contributions and explain negative coefficients.
\end{enumerate}
Record settings, numerical results, and observations.
\href{https://sosekie.github.io/Unibe-CVG-CV-Interactive-Labs/tutorial-09/#mixture}{Open the corresponding experiment.}\par
\Needspace{12\baselineskip}\section*{2. Covariance and a 2D input domain}
\noindent\textbf{Objective.} Read Gaussian geometry from covariance eigenpairs and distinguish scale from normalization.
\Needspace{9\baselineskip}\subsection*{2.1 Covariance - Principal Axes}
A 2D Gaussian has $\Sigma=\begin{bmatrix}0.20&0.16\\0.16&0.20\end{bmatrix}$.
The diagonal entries are equal.
\begin{enumerate}[label=(\alph*),nosep]
\item What shape does this covariance describe?
Justify your answer.
\item What happens to the usefulness of a rotation gradient as the two principal scales become equal?
\end{enumerate}
\Needspace{9\baselineskip}\subsection*{2.2 Gaussian Behavior - Normalization}
Consider one unnormalized 2D kernel $a\exp(-\frac12(x-\mu)^\top\Sigma^{-1}(x-\mu))$ with positive amplitude a > 0.
Double both principal scales while keeping a fixed.
\begin{enumerate}[label=(\alph*),nosep]
\item What happens to its peak value and its integral over the whole plane?
Justify your answer.
\item Does the fourfold integral increase imply fourfold image brightness after projecting and alpha-compositing a 3D splat?
\end{enumerate}
\Needspace{7\baselineskip}\noindent\textbf{Interactive experiment 2.}\begin{enumerate}[nosep]
\item Set one Gaussian’s scales to 0.60 and 0.20 and its rotation to 45${}^\circ$.
Inspect its covariance.
\item Make it isotropic and rotate again.
Explain which quantities stop changing and why.
\item Isolate one contribution and double both scales at fixed amplitude.
Predict the peak and full-plane integral; account for the finite plot and automatic display scaling.
\end{enumerate}
Record settings, numerical results, and observations.
\href{https://sosekie.github.io/Unibe-CVG-CV-Interactive-Labs/tutorial-09/#dimensions}{Open the corresponding experiment.}\par
\Needspace{12\baselineskip}\section*{3. Fit a 3D scalar function}
\noindent\textbf{Objective.} Separate approximation capacity, initialization, and optimization error.
\Needspace{9\baselineskip}\subsection*{3.1 Optimization - Gradient Flow}
The 3D model is $\hat{y}$(x) = $a\phi(x;\mu,\Sigma)$, fitted by mean squared error to sampled target values.
Consider the raw gradient at a = 0, with positive scales.
\begin{enumerate}[label=(\alph*),nosep]
\item Which parameters can receive a nonzero loss gradient at that instant?
Justify your answer.
\item Can Adam momentum move a mean when its current raw gradient is zero?
\end{enumerate}
\Needspace{9\baselineskip}\subsection*{3.2 Approximation - Capacity Vs Optimization}
One Gaussian fits the ring-shaped 3D target poorly after many updates.
Its slices fill the hole or miss much of the ring.
\begin{enumerate}[label=(\alph*),nosep]
\item What can be concluded from the Gaussian model itself, without assuming the optimizer found a global minimum?
Justify your answer.
\item If the displayed volume RMSE is small, could a visually important part of the ring still be poorly fitted?
\end{enumerate}
\Needspace{7\baselineskip}\noindent\textbf{Interactive experiment 3.}\begin{enumerate}[nosep]
\item Choose the Gaussian-bump target and N = 1.
Set amplitude to zero, record geometry, and apply one update.
Explain the raw-gradient effect.
\item Fit the ring target with N = 1 and N = 8 using comparable update budgets.
Inspect all three residual slices.
\item Rotate the inspection view without editing the fit.
Explain why the sampled scalar objective stays fixed.
\end{enumerate}
Record settings, numerical results, and observations.
\href{https://sosekie.github.io/Unibe-CVG-CV-Interactive-Labs/tutorial-09/#parameters}{Open the corresponding experiment.}\par
\Needspace{12\baselineskip}\section*{4. Project, composite, and expose ambiguity}
\noindent\textbf{Objective.} Relate visibility to gradients and distinguish rotation from translation.
\Needspace{9\baselineskip}\subsection*{4.1 Visibility - Photometric Gradients}
At one pixel, a foreground splat has $\alpha_1=0.98$ and a rear splat has $\alpha_2=0.6$.
Hold geometry and opacity fixed while differentiating the pixel colour with respect to the rear RGB.
\begin{enumerate}[label=(\alph*),nosep]
\item What is $\partial$C/$\partial$c${}_2$ for each RGB channel, and what does it imply for fitting?
Justify your answer.
\item Would directly adding the two coloured kernels give the same learning behavior?
\end{enumerate}
\Needspace{9\baselineskip}\subsection*{4.2 Multi-View Geometry - Ambiguity}
For an isolated Gaussian and a fixed pinhole camera, multiply its camera-relative mean and all spatial scales by k > 0.
Keep RGB and opacity fixed, and keep it in front of the near plane.
Its locally projected footprint is unchanged.
\begin{enumerate}[label=(\alph*),nosep]
\item Which extra observation generally constrains this depth-scale ambiguity?
Justify your answer.
\item Why can three nearby training views provide less useful depth information than three spread-out views?
\end{enumerate}
\Needspace{7\baselineskip}\noindent\textbf{Interactive experiment 4.}\begin{enumerate}[nosep]
\item Use three-splat overlap.
Sample one pixel, vary foreground opacity, and follow the rear colour weights.
\item Keep one Gaussian.
Record its projected axes, move it along the viewing ray with the provided scale change, and compare.
\item Change camera yaw to expose the 3D difference.
Here yaw orbits the camera centre; explain why pure rotation at a fixed centre is different.
\end{enumerate}
Record settings, numerical results, and observations.
\href{https://sosekie.github.io/Unibe-CVG-CV-Interactive-Labs/tutorial-09/#colour}{Open the corresponding experiment.}\par
\Needspace{12\baselineskip}\section*{5. Recover a scene from images}
\noindent\textbf{Objective.} Design controlled comparisons and distinguish image agreement from geometric evidence.
\Needspace{9\baselineskip}\subsection*{5.1 Experiment Design - Generalization}
A fit from one selected camera has low training MSE but a much larger reserved-view MSE.
You want to test whether limited view coverage is contributing to the failure.
\begin{enumerate}[label=(\alph*),nosep]
\item Which comparison provides the most relevant controlled evidence?
Justify your answer.
\item If you repeatedly choose budgets and stop times using frame 10, is it still an untouched final test?
\end{enumerate}
\Needspace{9\baselineskip}\subsection*{5.2 Density Control - Model Changes}
Consider one isolated coloured splat with opacity o.
Replace it by two coincident splats with identical shape and colour, each with opacity $o^\prime=1-\sqrt{1-o}$.
Let $\phi$ be their shared projected Gaussian value at a pixel.
\begin{enumerate}[label=(\alph*),nosep]
\item Does this opacity adjustment preserve the rendered image everywhere?
Justify your answer.
\item For an additive scalar field, would two coincident copies with amplitudes a/2 preserve the function?
\end{enumerate}
\Needspace{7\baselineskip}\noindent\textbf{Interactive experiment 5.}\begin{enumerate}[nosep]
\item Fit one training view with 24 Gaussians for 100 updates.
Record training and reserved-view MSE.
\item Restart with three spread-out views, using the same budget and update count.
Compare both errors and residual structure.
\item Record error and count, apply density refinement once, then fit 25 further updates.
Separate the immediate model change from subsequent learning.
\end{enumerate}
Record settings, numerical results, and observations.
\href{https://sosekie.github.io/Unibe-CVG-CV-Interactive-Labs/tutorial-09/#reconstruction}{Open the corresponding experiment.}\par
\Needspace{12\baselineskip}\section*{6. NeRF: a neural field along camera rays}
\noindent\textbf{Objective.} Distinguish density, segment opacity, quadrature, and neural-field capacity.
\Needspace{9\baselineskip}\subsection*{6.1 Volume Rendering - Optical Thickness}
A ray crosses a homogeneous slab of length $L=1$ with density $\sigma=2$ and constant RGB.
You approximate it with N equal intervals of length $\delta=L/N$, using $\alpha_i=1-\exp(-\sigma\delta)$.
\begin{enumerate}[label=(\alph*),nosep]
\item What happens to total transmittance if you change N from 10 to 100?
Justify your answer.
\item If a thin high-density feature lies entirely between the sample points, does a low sampled density prove the ray is empty?
\end{enumerate}
\Needspace{9\baselineskip}\subsection*{6.2 Neural Fields - Multi-View Consistency}
In a classical NeRF-style model, density is $\sigma_\theta(x)$, while RGB is $c_\theta(x,d)$.
A proposed variant also lets density depend freely on viewing direction: $\sigma_\theta(x,d)$.
\begin{enumerate}[label=(\alph*),nosep]
\item Why is position-only density a useful constraint for reconstructing a static scene?
Justify your answer.
\item Is increasing the number of samples per ray the NeRF equivalent of adding more Gaussian primitives?
\end{enumerate}
\Needspace{7\baselineskip}\noindent\textbf{Interactive experiment 6.}\begin{enumerate}[nosep]
\item At yaw 90${}^\circ$ and the centre ray, compare density with termination weights.
Explain the weak contribution of hidden orange density.
\item Compare 8 and 64 ray samples with fixed network weights.
Track interval length, alpha, and total transmittance.
\item Fit 50 updates and inspect the reserved 0${}^\circ$ view.
Compare errors and explain why additional ray samples do not add network parameters.
\end{enumerate}
Record settings, numerical results, and observations.
\href{https://sosekie.github.io/Unibe-CVG-CV-Interactive-Labs/tutorial-09/#nerf}{Open the corresponding experiment.}\par
\Needspace{12\baselineskip}\section*{7. Transfer to the real captured Garden}
\noindent\textbf{Objective.} Evaluate appearance and geometry claims at scene scale.
\Needspace{9\baselineskip}\subsection*{7.1 Appearance - Representation Choice}
Suppose a glossy leaf changes brightness with viewing direction in the captured photographs, but a reconstruction represents each Gaussian with one fixed RGB value.
Even with good spatial coverage, a highlight appears painted onto the surface.
\begin{enumerate}[label=(\alph*),nosep]
\item Which model change directly addresses this missing behavior?
Justify your answer.
\item Could a low photometric error prove that positions and covariances are correct when colour is very flexible?
\end{enumerate}
\Needspace{9\baselineskip}\subsection*{7.2 Transfer - Rendering Vs Geometry}
The garden looks sharp from several viewpoints.
A colleague wants to treat the visible Gaussian ellipsoids as a reliable surface model for collision detection.
\begin{enumerate}[label=(\alph*),nosep]
\item What does the rendering evidence justify?
Justify your answer.
\item If rendering at twice the resolution reveals smaller floaters, has the underlying geometry become worse?
\end{enumerate}
\Needspace{7\baselineskip}\noindent\textbf{Interactive experiment 7.}\begin{enumerate}[nosep]
\item Load the Garden when network and WebGL are available.
It streams approximately 187 MB and renders Gaussian Splatting.
\item Inspect thin structures and viewpoints away from the comfortable viewing region.
Record two artifacts and competing explanations.
\item Relate observations to coverage, representation, or appearance; identify evidence needed before using this renderer as collision geometry.
\end{enumerate}
Record settings, numerical results, and observations.
\href{https://sosekie.github.io/Unibe-CVG-CV-Interactive-Labs/tutorial-09/#highres}{Open the corresponding experiment.}\par
\end{document}
